% =====================================================================
%  Fiche pédagogique — Physique, chapitre 8 : L’énergie potentielle électrostatique (sciences mathématiques, 10 h, cinq séances)
% =====================================================================
\newcommand{\Matiere}{Physique}
\newcommand{\Partie}{Partie 2 — Électricité}
\newcommand{\NumChap}{8}
\newcommand{\TitreChap}{L’énergie potentielle\\ électrostatique}
\newcommand{\TitreCourt}{L’énergie potentielle électrostatique}
\newcommand{\Niveau}{1\textsuperscript{re} année du baccalauréat}
\newcommand{\Duree}{10 h}
\newcommand{\Banniere}{opener_phys8.jpg}
\newcommand{\BanniereDecalage}{1.3cm}
\newcommand{\Epe}{E_{\mathrm{pe}}}
\newcommand{\Ecin}{E_{\mathrm{C}}}
\newcommand{\Epp}{E_{\mathrm{pp}}}
\newcommand{\vect}[1]{\overrightarrow{#1}}
\newcommand{\W}[2]{W_{#1\rightarrow #2}(\vect{F})}

\begin{document}
\bandeau

\noindent
\begin{minipage}[t]{0.235\linewidth}\infocarte{fi_toque}{cPurple}{Niveau}{\Niveau\\ Sciences mathématiques (option française)}\end{minipage}\hfill
\begin{minipage}[t]{0.235\linewidth}\infocarte{fi_horloge}{cAmber}{Durée}{10 heures — 5 séances\\ 8 h cours + 2 h exercices}\end{minipage}\hfill
\begin{minipage}[t]{0.235\linewidth}\infocarte{fi_prof}{cCyan}{Professeur}{\rule{0pt}{3.6mm}{\color{cInk2!60}\dotfill}\\ Classe : {\color{cInk2!60}\dotfill}}\end{minipage}\hfill
\begin{minipage}[t]{0.235\linewidth}\infocarte{fi_chemin}{cInk}{Programme}{Physique · Électricité · § 1 (SM)\\ Énergie potentielle électrostatique}\end{minipage}

% =====================================================================
\bloc{1}{fi_cible}{Cadre pédagogique de la séquence}

\begin{xltabular}{\linewidth}{@{}C{3.5cm}I Y@{}}
\ligne{fi_cible}{cRed}{Compétence\\ spécifique}{%
\textbf{Électricité :} \textbf{interpréter les transferts d’énergie} et les phénomènes de \textbf{dissipation et de conservation} dans des circuits électriques, dans diverses situations ; \textbf{résoudre une situation-problème} relative au bilan énergétique dans des circuits électriques, expérimentalement ou par une étude analytique {\color{cInk2}(compétence spécifique de la physique, 1\textsuperscript{re} année du baccalauréat)}.
\pt{\textbf{Contribution de ce chapitre :} premier chapitre de la partie Électricité (propre à la filière sciences mathématiques) : \textbf{électrisation} et \textbf{loi de Coulomb} ; \textbf{champ électrostatique} (charge ponctuelle, superposition, lignes de champ, champ uniforme) ; \textbf{travail de la force électrostatique}, \textbf{potentiel} et différence de potentiel ; \textbf{énergie potentielle électrostatique} et \textbf{conservation de l’énergie totale} d’une particule chargée. Situations-problèmes : atome d’hydrogène, électron entre deux plaques, canon à électrons, pendule et ion dans un condensateur plan.}}
\ligne{fi_reseau}{cTeal}{Compétences\\ transversales}{%
\begin{minipage}[t]{0.493\linewidth}\RaggedRight
\pt{\textbf{Démarche scientifique} — maîtriser les éléments de la démarche scientifique dans toutes ses dimensions.\hfill\break{\color{cTeal}\itshape Ici :} expériences d’électrisation et interprétation par un modèle (électrons) ; mesures avec des charges témoins ; spectres électrostatiques.}
\pt{\textbf{Communication} — lire, écrire, écouter, dialoguer ; maîtriser la langue écrite et le langage mathématique ; utiliser les symboles ; respecter l’éthique du dialogue.\hfill\break{\color{cTeal}\itshape Ici :} notations vectorielles ($\vect{F}_{A/B}$, $\vect{E}$, $\vect{AB}$), schémas de forces et de champs, rédaction de démonstrations.}
\pt{\textbf{Technologies de l’information} — utiliser les TIC pour s’informer, traiter des données, éprouver des hypothèses, construire des modèles.\hfill\break{\color{cTeal}\itshape Ici :} animations (électrisation, pendules, spectres, plaques chargées, canon à électrons).}
\end{minipage}\hfill
\begin{minipage}[t]{0.483\linewidth}\RaggedRight
\pt{\textbf{Méthodologie} — acquérir une méthodologie de recherche, de travail et d’autoformation ; réaliser des tâches complémentaires.\hfill\break{\color{cTeal}\itshape Ici :} méthode de superposition des champs ; analogie avec le travail du poids et l’énergie potentielle de pesanteur.}
\pt{\textbf{Culture intégrée} — s’intégrer dans son environnement, citoyenneté active, positions positives face aux questions d’environnement, de santé, de prévention et de consommation.\hfill\break{\color{cTeal}\itshape Ici :} décharges électrostatiques et sécurité (éclairs, paratonnerre, stations-service) ; tubes cathodiques et accélérateurs de particules.}
\end{minipage}}
\ligne{fi_drapeau}{cBlue}{Objectifs\\ d’apprentissage}{%
{\color{cInk2}\itshape À la fin de la séquence, l’apprenant(e) est capable de :}\par\vspace{.6mm}
\begin{minipage}[t]{0.493\linewidth}\RaggedRight
\pt{\textbf{O1} — interpréter l’\textbf{électrisation} (frottement, contact, influence) par des transferts ou des déplacements d’électrons ;}
\pt{\textbf{O2} — connaître et appliquer la \textbf{loi de Coulomb} (forme vectorielle, intensité) ;}
\pt{\textbf{O3} — connaître le \textbf{champ électrostatique}, utiliser $\vect{F} = q\cdot\vect{E}$ et le champ d’une charge ponctuelle ; appliquer la superposition ;}
\pt{\textbf{O4} — reconnaître les \textbf{lignes de champ} (charge ponctuelle, deux charges) et le \textbf{champ uniforme} entre deux plaques ;}
\end{minipage}\hfill
\begin{minipage}[t]{0.483\linewidth}\RaggedRight
\pt{\textbf{O5} — établir et exploiter $\W{A}{B} = q\cdot(V_A - V_B)$, $V_A - V_B = \vect{E}\cdot\vect{AB}$ et $E = U/d$ ; reconnaître les équipotentielles ;}
\pt{\textbf{O6} — connaître et exploiter $\Epe = q\cdot V + C$ et $\Delta\Epe = -\W{A}{B}$ ;}
\pt{\textbf{O7} — appliquer la \textbf{conservation de l’énergie totale} d’une particule chargée ; utiliser l’\textbf{électron-volt}.}
\end{minipage}
\par\vspace{1mm}{\footnotesize\color{cInk2}\textit{Connaissances et habiletés du programme :} connaître et appliquer la loi de Coulomb ; connaître le champ électrostatique ; connaître et appliquer la relation $\vect{E} = \vect{F}/q$ ; reconnaître une ligne de champ ; connaître la forme des lignes de champ pour une charge ponctuelle et pour deux charges ponctuelles ; connaître et exploiter $W = q(V_A - V_B)$, où $V_A - V_B$ est la différence de potentiel et $V$ le potentiel électrique en un point du champ ; connaître et exploiter $\Epe = q\cdot V + C$, énergie potentielle électrostatique en un point du champ.}}
\ligne{fi_marches}{cGreen}{Prérequis}{%
\begin{minipage}[t]{0.493\linewidth}\RaggedRight
{\sffamily\bfseries\footnotesize\color{cGreen}Acquis des chapitres 2, 3 et 4 et du tronc commun}
\pt{travail d’une force constante $W = \vect{F}\cdot\vect{AB}$ (chapitre 2) ; théorème de l’énergie cinétique (chapitre 3) ;}
\pt{énergie potentielle de pesanteur $\Epp = m\cdot g\cdot z + C$, énergie mécanique (chapitre 4) ;}
\pt{structure de l’atome, charge élémentaire ; loi de gravitation universelle ; tension $U_{AB} = V_A - V_B$.}
\end{minipage}\hfill
\begin{minipage}[t]{0.483\linewidth}\RaggedRight
{\sffamily\bfseries\footnotesize\color{cGreen}Outils à réactiver}
\pt{vecteurs : somme, projection, produit scalaire en coordonnées ;}
\pt{trigonométrie ($\cos 60^\circ$, $\tan\alpha$), théorème de Pythagore ;}
\pt{puissances de 10, préfixes (n, µ), conversions cm $\to$ m.}
\end{minipage}}
\ligne{fi_boite}{cOrange}{Matériel\\ didactique}{%
\begin{minipage}[t]{0.493\linewidth}\RaggedRight
{\sffamily\bfseries\footnotesize\color{cOrange}Expériences de cours}
\pt{règle en plastique, tissu de laine, petits morceaux de papier ; bâtons de verre et d’ébonite, soie, fourrure, étrier et potence ;}
\pt{pendules électrostatiques (boules légères métallisées, fils isolants) ;}
\pt{générateur haute tension (ou machine électrostatique), sphère conductrice sur pied isolant ;}
\pt{cuve transparente, huile de paraffine, semoule, électrodes ponctuelles et plaques ; deux plaques métalliques parallèles sur supports isolants ;}
\pt{oscilloscope ou tube à faisceau d’électrons, s’il est disponible (sinon l’animation du canon à électrons).}
\end{minipage}\hfill
\begin{minipage}[t]{0.483\linewidth}\RaggedRight
{\sffamily\bfseries\footnotesize\color{cOrange}Supports}
\pt{vidéoprojecteur et ordinateur : \textbf{présentation animée} du chapitre (42 diapositives) ;}
\pt{\textbf{cahier} de l’apprenant : fiche d’activités collée, réponses aux questions, « À retenir », applications et exercices ;}
\pt{\textbf{fiche d’activités} imprimée (deux pages : activités 1 et 2, puis 3 et 4, à coller dans le cahier) ;}
\pt{tableau, feutres ou craies de couleur, calculatrices.}
\pt{\textbf{Sécurité :} la haute tension est manipulée par le professeur seul, générateur à courant limité.}
\end{minipage}}
\ligne{fi_livre}{cInk2}{Références}{%
\pt{Ministère de l’Éducation nationale — \textit{Orientations pédagogiques générales et programmes spécifiques à l’enseignement de la physique et de la chimie au cycle secondaire qualifiant}, 2007 : programme de la 1\textsuperscript{re} année du baccalauréat — physique, Électricité, § 1 « Énergie potentielle électrostatique » (propre à la filière sciences mathématiques, 8 h de cours et 2 h d’exercices) : 1.1 le champ électrostatique (interaction électrostatique, loi de Coulomb ; champ créé par une charge ponctuelle : définition, direction, sens ; exemples de lignes de champ ; superposition de deux champs ; champ uniforme) ; 1.2 énergie potentielle d’une charge dans un champ électrostatique uniforme (travail de la force électrostatique ; potentiel et différence de potentiel ; surfaces équipotentielles ; relation entre énergie potentielle et travail ; énergie totale d’une particule chargée soumise à la seule force électrostatique et sa conservation).}
\pt{Activités proposées, toutes traitées dans la séquence : expériences d’électrisation de la matière (frottement, contact, influence) ; mise en évidence expérimentale du champ électrostatique ; visualisation des lignes de champ avec l’huile de paraffine et la semoule ; réalisation d’un champ uniforme avec deux plaques métalliques parallèles ; établissement du travail de la force électrostatique et lien avec la différence de potentiel et l’énergie potentielle électrostatique.}
\pt{Supports : présentation animée et fiche d’activités ; manuels scolaires agréés de la 1\textsuperscript{re} année du baccalauréat (sciences mathématiques).}}
\end{xltabular}

% =====================================================================
\needspace{9cm}
\bloc{2}{fi_chemin}{Déroulement de la séquence (cinq séances)}

\seance{1}{2 h}{électrisation et loi de Coulomb}
\frise[120]{cCyan/10/Mise en route, cTeal/103/Partie 1 · Électrisation et loi de Coulomb, cGreen/7/Bilan}
\seance{2}{2 h}{le champ électrostatique}
\frise[120]{cCyan/8/Rappels, cBlue/107/Partie 2 · Champ électrostatique, cGreen/5/Bilan}
\seance{3}{2 h}{travail de la force électrostatique et potentiel}
\frise[120]{cCyan/18/Rappels et exercice 1, cPurple/97/Partie 3 · Travail et potentiel, cGreen/5/Bilan}
\seance{4}{2 h}{énergie potentielle électrostatique}
\frise[120]{cCyan/8/Rappels, cAmber/97/Partie 4 · Énergie potentielle, cOrange/10/Vrai ou faux, cGreen/5/Bilan}
\seance{5}{2 h}{exercices de synthèse et évaluation}
\frise[120]{cCyan/5/Rappels, cOrange/100/Exercices 2 à 4, cGreen/15/L’essentiel et bilan}
\par\vspace{4mm}

\begin{xltabular}{\linewidth}{@{}L{3.75cm}I >{\small\RaggedRight\arraybackslash}XI >{\small\RaggedRight\arraybackslash}XI L{4.6cm}I C{1.45cm}@{}}
\enteteDeroulement
\endfirsthead
\enteteDeroulement
\endhead
% =====================================================================
\partie{cInk}{Séance 1 · Électrisation et loi de Coulomb (2 h)}{0 h 00 $\to$ 2 h 00}
\etape{cCyan}{\axe{cInk2}{ic_rappel}{Mise en route}{Accueil, plan et prérequis}{\diapo{1} \support{Fiche}}}
{\pt{Distribue la fiche d’activités (deux pages), que chacun colle dans son cahier ; annonce le plan des cinq séances.}
 \pt{Questions flash : « de quoi est formé un atome ? », « charge d’un électron ? », « expression du travail d’une force constante ? », « loi de gravitation entre deux masses ? ».}}
{\pt{Répond oralement : noyau (protons, neutrons) et électrons ; $-e = \qty{-1,6e-19}{C}$ ; $W = \vect{F}\cdot\vect{AB}$ ; $F = G\cdot m_A\cdot m_B/d^2$.}
 \pt{Colle la fiche d’activités dans le cahier.}}
{\diag Questions flash.\par
 \textit{Indicateurs :} structure de l’atome ; produit scalaire.\par\vspace{.8mm}
 \regul Rappel au tableau du produit scalaire en coordonnées.}
{10}{0 h 10}
% ---------------------------------------------------------------------
\partie{cTeal}{Partie 1 · Électrisation et loi de Coulomb}{103 min · objectifs O1 et O2}
\etape{cTeal}{\axe{cAmber}{ic_situation}{Situation de départ}{La règle frottée, la petite décharge}{\diapo{2}}}
{\pt{Frotte une règle en plastique avec un tissu de laine et l’approche de petits morceaux de papier ; projette l’animation (diapo 2).}
 \pt{Pose le problème : que gagne un corps frotté, comment l’expliquer, quelle force agit ?}}
{\pt{Observe ; propose : « la règle s’est chargée », « elle attire comme un aimant »…}}
{\diag Représentations initiales (confusion avec le magnétisme).}
{8}{0 h 18}
\etape{cTeal}{\axe{cTeal}{ic_activite}{Activité 1}{Deux sortes de charges ; contact et influence (docs 1 à 3)}{\diapo{3--6} \support{Fiche · Act. 1} \support{Cahier}}}
{\pt{Réalise les expériences des bâtons suspendus (verre-verre, ébonite-ébonite, verre-ébonite), puis du pendule électrostatique par contact et par influence (diapos 3 et 4).}
 \pt{Présente l’information 1 ; fait traiter les questions 1 et 2 en binômes ; corrige au clic (diapo 6).}}
{\pt{Observe, décrit, rédige ses réponses dans le cahier.}
 \pt{Conclut à deux sortes de charges et à la règle « même signe : répulsion » ; explique le frottement et le contact par un transfert d’électrons, l’influence par un déplacement des électrons libres (boule neutre, attraction en $1/d^2$ plus forte).}}
{\form Observation des binômes, correction collective.\par
 \textit{Indicateurs :} seuls les électrons se déplacent ; neutralité de la boule influencée.\par\vspace{.8mm}
 \regul Erreur fréquente : « des charges positives passent » ; revenir à l’information 1.}
{45}{1 h 03}
\etape{cTeal}{\axe{cTeal}{ic_questions}{Activité 1 (suite)}{La loi de Coulomb (doc 4)}{\diapo{5, 7} \support{Cahier}}}
{\pt{Présente la loi de Coulomb (diapo 5) ; fait préciser les caractéristiques des deux forces et les représenter (questions 3 et 4, diapo 7).}}
{\pt{Donne point d’application, droite d’action, sens et intensité ; représente les forces répulsives puis attractives.}}
{\form Schémas au tableau par deux apprenants.\par
 \textit{Indicateurs :} $\vect{F}_{A/B}$ appliquée en $B$ ; $\vect{F}_{A/B} = -\vect{F}_{B/A}$.}
{20}{1 h 23}
\etape{cTeal}{\axe{cRed}{ic_retenir}{À retenir}{Électrisation ; loi de Coulomb}{\diapo{8--9} \support{Cahier}}}
{\pt{Projette les diapos 8 et 9 et laisse le temps de recopier.}}
{\pt{Recopie l’« À retenir » dans son cahier.}}
{\form Question flash : « la distance entre deux charges triple : que devient $F$ ? » $\to$ divisée par 9.}
{12}{1 h 35}
\etape{cTeal}{\axe{cPurple}{ic_application}{Application}{L’atome d’hydrogène}{\diapo{10} \support{Cahier}}}
{\pt{Projette l’énoncé (diapo 10) ; travail individuel ; correction au clic.}}
{\pt{Calcule $F_e \approx \qty{8,2e-8}{N}$, $F_G \approx \qty{3,7e-47}{N}$, $F_e/F_G \approx \num{2,2e39}$ ; conclut que la gravitation est négligeable à l’échelle atomique.}}
{\somm Exercice individuel.\par
 \textit{Indicateur :} puissances de 10 correctes.}
{18}{1 h 53}
\etape{cGreen}{\axe{cGreen}{ic_reponses}{Bilan}{Synthèse et travail à la maison}{\support{Cahier}}}
{\pt{Fait reformuler la règle d’attraction et la loi de Coulomb ; donne le travail à la maison.}}
{\pt{Reformule ; note le travail.}}
{\form Reformulation orale.}
{7}{2 h 00}
% =====================================================================
\partie{cInk}{Séance 2 · Le champ électrostatique (2 h)}{0 h 00 $\to$ 2 h 00}
\etape{cCyan}{\axe{cInk2}{ic_rappel}{Rappels}{Loi de Coulomb}{\diapo{9}}}
{\pt{Questions flash : « expression de $F$ ? », « valeur de $K$ ? », « point d’application de $\vect{F}_{A/B}$ ? ».}}
{\pt{Répond : $F = K\cdot|q_A|\cdot|q_B|/AB^2$ ; \qty{9e9}{N.m^2.C^{-2}} ; en $B$.}}
{\diag Questions flash.}
{8}{0 h 08}
\partie{cBlue}{Partie 2 · Le champ électrostatique}{107 min · objectifs O3 et O4}
\etape{cBlue}{\axe{cAmber}{ic_situation}{Situation de départ}{Une action à distance}{\diapo{11}}}
{\pt{Approche un pendule chargé d’une sphère chargée ; projette l’animation (diapo 11) ; pose le problème : comment caractériser la propriété de l’espace autour d’une charge ?}}
{\pt{Observe la déviation à distance ; propose « une zone d’influence ».}}
{\diag Hypothèses.}
{7}{0 h 15}
\etape{cBlue}{\axe{cTeal}{ic_activite}{Activité 2}{Pendules déviés, charges témoins (docs 5 et 6)}{\diapo{12, 14} \support{Fiche · Act. 2} \support{Cahier}}}
{\pt{Déplace le pendule chargé autour de la sphère (doc 5) ; exploite le tableau des charges témoins (doc 6, diapo 12) et l’information 2 ; fait traiter les questions 1 et 2 ; corrige (diapo 14).}}
{\pt{Constate la déviation qui diminue avec la distance ; calcule $F_i/|q_i| = \qty{5,0e3}{N.C^{-1}}$ pour les trois charges ; en déduit $\vect{E}$ en $P$ et compare les sens de $\vect{F}$ et de $\vect{E}$.}}
{\form Calculs vérifiés en binômes.\par
 \textit{Indicateurs :} rapport indépendant de $q_i$ ; sens de $\vect{F}$ pour $q < 0$.}
{28}{0 h 43}
\etape{cBlue}{\axe{cTeal}{ic_questions}{Activité 2 (suite)}{Charge ponctuelle, superposition}{\diapo{15} \support{Cahier}}}
{\pt{Guide l’établissement de $\vect{E} = K\cdot\dfrac{q}{r^2}\,\vect{u}$ à partir de la loi de Coulomb ; fait vérifier en $P$ et en $P'$ ; questions 3 et 4 (diapo 15).}}
{\pt{Établit $\vect{E}$ et $E$ ; trouve \num{5,0e3} et \qty{1,25e3}{N.C^{-1}} ; distingue champ centrifuge et centripète ; écrit $\vect{E} = \sum\vect{E}_i$.}}
{\form Correction collective.\par
 \textit{Indicateur :} le vecteur $\vect{u}$ va de la charge vers le point.}
{20}{1 h 03}
\etape{cBlue}{\axe{cTeal}{ic_activite}{Activité 2 (suite)}{Spectres électrostatiques (doc 7)}{\diapo{13, 16} \support{Fiche · Act. 2}}}
{\pt{Réalise les spectres (huile de paraffine, semoule, haute tension) avec une électrode, deux électrodes, deux plaques ; présente l’information 3 ; question 5 (diapo 16).}}
{\pt{Décrit les lignes radiales, courbes, parallèles ; justifie que le champ est uniforme entre les plaques, de la plaque positive vers la négative.}}
{\form Description orale guidée.\par
 \textit{Indicateur :} lignes orientées des charges positives vers les négatives.}
{20}{1 h 23}
\etape{cBlue}{\axe{cRed}{ic_retenir}{À retenir}{$\vect{F} = q\cdot\vect{E}$ ; lignes de champ}{\diapo{17--18} \support{Cahier}}}
{\pt{Projette les diapos 17 et 18 et laisse le temps de recopier.}}
{\pt{Recopie l’« À retenir » dans son cahier.}}
{\form Question flash : « sens de $\vect{E}$ autour d’une charge négative ? » $\to$ vers la charge.}
{14}{1 h 37}
\etape{cBlue}{\axe{cOrange}{ic_methode}{Méthode}{Champ créé par plusieurs charges}{\diapo{19} \support{Cahier}}}
{\pt{Présente la méthode en quatre étapes (diapo 19) sur l’exemple de deux charges opposées.}}
{\pt{Note la méthode dans le cahier.}}
{\form Questions de compréhension.}
{8}{1 h 45}
\etape{cBlue}{\axe{cPurple}{ic_application}{Application}{Champ au milieu de deux charges opposées}{\diapo{20} \support{Cahier}}}
{\pt{Projette l’énoncé (diapo 20) ; travail individuel ; correction au clic.}}
{\pt{Représente $\vect{E}_A$ et $\vect{E}_B$, de même sens ; calcule $E(O) = \qty{1,44e7}{V.m^{-1}}$, de $A$ vers $B$.}}
{\somm Exercice individuel.\par
 \textit{Indicateur :} sens de $\vect{E}_B$ (vers la charge négative).}
{10}{1 h 55}
\etape{cGreen}{\axe{cGreen}{ic_reponses}{Bilan}{Travail à la maison}{\diapo{34}}}
{\pt{Donne l’exercice 1 (énoncé projeté, diapo 34) à résoudre dans le cahier.}}
{\pt{Recopie l’énoncé.}}
{\form —}
{5}{2 h 00}
% =====================================================================
\partie{cInk}{Séance 3 · Travail de la force électrostatique et potentiel (2 h)}{0 h 00 $\to$ 2 h 00}
\etape{cCyan}{\axe{cInk2}{ic_rappel}{Rappels}{Champ électrostatique}{\diapo{17}}}
{\pt{Questions flash : « relation entre $\vect{F}$ et $\vect{E}$ ? », « champ entre deux plaques ? », « travail d’une force constante ? ».}}
{\pt{Répond : $\vect{F} = q\cdot\vect{E}$ ; uniforme, de $P$ vers $N$ ; $W = \vect{F}\cdot\vect{AB}$.}}
{\diag Questions flash.}
{8}{0 h 08}
\etape{cCyan}{\axe{cOrange}{ic_exercices}{Exercice 1}{Correction du travail à la maison}{\diapo{35} \support{Cahier}}}
{\pt{Corrige l’exercice 1 au clic (diapo 35) ; commente la courbe $E = f(x)$.}}
{\pt{Compare avec sa solution : $E_p = \qty{14,4}{N.C^{-1}}$, $q \approx \qty{1,11e-13}{C}$, décroissance en $1/x^2$.}}
{\form Autocorrection.\par\vspace{.8mm}
 \regul Conversions µm et cm $\to$ m.}
{10}{0 h 18}
\partie{cPurple}{Partie 3 · Travail de la force électrostatique et potentiel électrique}{97 min · objectif O5}
\etape{cPurple}{\axe{cAmber}{ic_situation}{Situation de départ}{Analogie avec le travail du poids}{\diapo{21}}}
{\pt{Rappelle $W(\vect{P}) = m\cdot g\cdot(z_A - z_B)$ ; projette l’animation des plaques (diapo 21) ; pose le problème.}}
{\pt{Propose que le travail dépende « de la position » ou « de la tension ».}}
{\diag Hypothèses.}
{7}{0 h 25}
\etape{cPurple}{\axe{cTeal}{ic_activite}{Activité 3}{Boule chargée entre deux plaques (doc 8)}{\diapo{22--23} \support{Fiche · Act. 3} \support{Cahier}}}
{\pt{Montre la déviation d’un pendule chargé entre deux plaques reliées à la haute tension (doc 8) ; présente le rappel et l’information 4 ; questions 1 et 2 en binômes ; corrige (diapo 23).}}
{\pt{Calcule $\W{A}{B} = q\cdot E\cdot(x_A - x_B)$, indépendant du chemin (force conservative) ; écrit $V_A - V_B = E\cdot x_A - E\cdot x_B$ et $\W{A}{B} = q\cdot(V_A - V_B)$.}}
{\form Observation des binômes.\par
 \textit{Indicateurs :} produit scalaire avec $\vect{E} = -E\,\vect{i}$ ; signe.\par\vspace{.8mm}
 \regul Refaire le calcul pour un déplacement parallèle aux plaques.}
{35}{1 h 00}
\etape{cPurple}{\axe{cTeal}{ic_questions}{Activité 3 (suite)}{Potentiels décroissants, équipotentielles, $E = U/d$}{\diapo{24} \support{Cahier}}}
{\pt{Fait traiter les questions 3 à 5 ; corrige au clic (diapo 24).}}
{\pt{Trouve $V_P - V_N = E\cdot d > 0$ ; montre que $\vect{E}$ va vers les potentiels décroissants, que les équipotentielles sont des plans parallèles aux plaques ; écrit $E = U_{PN}/d$.}}
{\form Correction collective.\par
 \textit{Indicateur :} $\vect{E}\cdot\vect{AC} = 0 \Rightarrow \vect{AC} \perp \vect{E}$.}
{25}{1 h 25}
\etape{cPurple}{\axe{cRed}{ic_retenir}{À retenir}{$W = q\cdot(V_A - V_B)$ ; $E = U/d$}{\diapo{25} \support{Cahier}}}
{\pt{Projette la diapo 25 et laisse le temps de recopier.}}
{\pt{Recopie l’« À retenir » dans son cahier.}}
{\form Question flash : « travail sur une équipotentielle ? » $\to$ nul.}
{15}{1 h 40}
\etape{cPurple}{\axe{cPurple}{ic_application}{Application}{Un électron entre deux plaques}{\diapo{26} \support{Cahier}}}
{\pt{Projette l’énoncé (diapo 26) ; travail individuel ; correction au clic.}}
{\pt{Calcule $U_{PN} = \qty{3,0e3}{V}$ et $\W{N}{P} = \qty{4,8e-16}{J}$, travail moteur.}}
{\somm Exercice individuel.\par
 \textit{Indicateur :} signe de $q = -e$ et ordre $V_N - V_P$.}
{15}{1 h 55}
\etape{cGreen}{\axe{cGreen}{ic_reponses}{Bilan}{Synthèse}{\support{Cahier}}}
{\pt{Fait comparer $W(\vect{P}) = m\cdot g\cdot(z_A - z_B)$ et $W(\vect{F}) = q\cdot(V_A - V_B)$.}}
{\pt{Établit l’analogie $m\cdot g \leftrightarrow q$, $z \leftrightarrow V$.}}
{\form Reformulation orale.}
{5}{2 h 00}
% =====================================================================
\partie{cInk}{Séance 4 · Énergie potentielle électrostatique (2 h)}{0 h 00 $\to$ 2 h 00}
\etape{cCyan}{\axe{cInk2}{ic_rappel}{Rappels}{Travail et potentiel ; théorème de l’énergie cinétique}{\diapo{25}}}
{\pt{Questions flash : « $\W{A}{B}$ ? », « énoncé du théorème de l’énergie cinétique ? », « expression de $\Epp$ ? ».}}
{\pt{Répond : $q\cdot(V_A - V_B)$ ; $\Delta\Ecin = \sum W$ ; $m\cdot g\cdot z + C$.}}
{\diag Questions flash.}
{8}{0 h 08}
\partie{cAmber}{Partie 4 · Énergie potentielle électrostatique et conservation de l’énergie}{97 min · objectifs O6 et O7}
\etape{cAmber}{\axe{cAmber}{ic_situation}{Situation de départ}{Le canon à électrons}{\diapo{27}}}
{\pt{Projette l’animation du canon à électrons (diapo 27) ou montre le spot d’un oscilloscope ; pose le problème.}}
{\pt{Propose : l’énergie cinétique vient « de la tension », « du champ ».}}
{\diag Hypothèses.}
{7}{0 h 15}
\etape{cAmber}{\axe{cTeal}{ic_activite}{Activité 4}{Énergie potentielle et travail (doc 9)}{\diapo{28--29} \support{Fiche · Act. 4} \support{Cahier}}}
{\pt{Présente le doc 9 et l’information 5 (diapo 28) ; questions 1 et 2 en binômes ; corrige (diapo 29).}}
{\pt{Donne l’unité (J) ; détermine $C$ avec la référence ; établit $\Delta\Epe = q\cdot(V_B - V_A) = -\W{A}{B}$.}}
{\form Observation des binômes.\par
 \textit{Indicateur :} analogie avec $\Delta\Epp = -W(\vect{P})$.}
{30}{0 h 45}
\etape{cAmber}{\axe{cTeal}{ic_questions}{Activité 4 (suite)}{Conservation de l’énergie, électron-volt}{\diapo{30} \support{Cahier}}}
{\pt{Fait démontrer la conservation de $E = \Ecin + \Epe$ (question 3) et calculer l’électron-volt (question 4) ; corrige (diapo 30).}}
{\pt{Démontre $\Ecin(B) + \Epe(B) = \Ecin(A) + \Epe(A)$ ; trouve $\qty{1}{eV} = \qty{1,6e-19}{J}$ et les multiples.}}
{\form Correction collective.\par
 \textit{Indicateur :} poids négligeable justifié par les ordres de grandeur.}
{25}{1 h 10}
\etape{cAmber}{\axe{cRed}{ic_retenir}{À retenir}{$\Epe = q\cdot V + C$ ; énergie totale}{\diapo{31} \support{Cahier}}}
{\pt{Projette la diapo 31 et laisse le temps de recopier.}}
{\pt{Recopie l’« À retenir » dans son cahier.}}
{\form Question flash : « énergie gagnée par un proton accéléré sous \qty{500}{V} ? » $\to$ \qty{500}{eV}.}
{15}{1 h 25}
\etape{cAmber}{\axe{cPurple}{ic_application}{Application}{Le canon à électrons}{\diapo{32} \support{Cahier}}}
{\pt{Projette l’énoncé (diapo 32) ; travail individuel ; correction au clic.}}
{\pt{Calcule $\Delta\Epe = \qty{-3000}{eV} = \qty{-4,8e-16}{J}$, $\Ecin(A) = \qty{4,8e-16}{J}$ et $v_A \approx \qty{3,25e7}{m.s^{-1}}$.}}
{\somm Exercice individuel.\par
 \textit{Indicateur :} conversion eV $\to$ J ; racine carrée.}
{20}{1 h 45}
\etape{cOrange}{\axe{cOrange}{ic_questions}{Vrai ou faux ?}{Je vérifie mes connaissances}{\diapo{33}}}
{\pt{Projette le vrai/faux interactif (diapo 33) ; interroge à main levée.}}
{\pt{Répond et corrige les affirmations fausses à l’oral.}}
{\form Vrai/faux à main levée.\par
 \textit{Indicateur :} au moins 5 réponses justes sur 7.\par\vspace{.8mm}
 \regul Reprendre la notion correspondante pour chaque erreur fréquente.}
{10}{1 h 55}
\etape{cGreen}{\axe{cGreen}{ic_reponses}{Bilan}{Travail à la maison}{\diapo{36}}}
{\pt{Donne l’exercice 2 (énoncé projeté, diapo 36) à préparer dans le cahier.}}
{\pt{Recopie l’énoncé.}}
{\form —}
{5}{2 h 00}
% =====================================================================
\partie{cInk}{Séance 5 · Exercices de synthèse et évaluation (2 h)}{0 h 00 $\to$ 2 h 00 · objectifs O2 à O7}
\etape{cCyan}{\axe{cInk2}{ic_rappel}{Rappels}{Les relations du chapitre}{\diapo{42}}}
{\pt{Fait citer les relations du chapitre à partir de la diapo « L’essentiel » (masquée au départ).}}
{\pt{Cite les relations et leurs unités.}}
{\diag Questions flash.}
{5}{0 h 05}
\etape{cOrange}{\axe{cOrange}{ic_exercices}{Exercice 2}{Superposition de champs}{\diapo{36--37} \support{Cahier}}}
{\pt{Fait présenter au tableau la solution préparée ; corrige au clic (diapo 37) en insistant sur les projections.}}
{\pt{Trouve $E(M) = \dfrac{K\cdot q}{4a^2}$, parallèle à $(AB)$, et $E(D) = \left(\sqrt2 - \tfrac12\right)\dfrac{K\cdot q}{a^2} \approx 0{,}914\,\dfrac{K\cdot q}{a^2}$.}}
{\form Présentation au tableau.\par
 \textit{Indicateurs :} sens de chaque $\vect{E}_i$ ; distances ($a\sqrt2$).\par\vspace{.8mm}
 \regul Schéma vectoriel repris pas à pas pour les apprenants en difficulté.}
{30}{0 h 35}
\etape{cOrange}{\axe{cOrange}{ic_exercices}{Exercice 3}{Une particule $\alpha$ dans un champ uniforme}{\diapo{38--39} \support{Cahier}}}
{\pt{Projette l’énoncé (diapo 38) ; travail en binômes ; correction au clic (diapo 39).}}
{\pt{Calcule $\W{A}{B} = \qty{6,4e-18}{J}$, $\Epe(B) = \qty{-6,4e-18}{J}$, $\Ecin = \qty{2,0}{MeV} = \qty{3,2e-13}{J}$.}}
{\form Correction collective.\par
 \textit{Indicateur :} seule la coordonnée $x$ intervient ($\vect{E} = E\,\vect{i}$).}
{25}{1 h 00}
\etape{cOrange}{\axe{cOrange}{ic_exercices}{Exercice 4}{Pendule et ion entre deux plaques}{\diapo{40--41} \support{Cahier}}}
{\pt{Projette l’énoncé (diapo 40) ; travail individuel noté ; ramasse quelques cahiers ; corrige au clic (diapo 41).}}
{\pt{Calcule $E = \qty{5,0e3}{V.m^{-1}}$, $\alpha \approx 5{,}7^\circ$, $V_M = \qty{250}{V}$, $\Epe(M) = \qty{1,0e-5}{J}$, $v_N \approx \qty{8,30e4}{m.s^{-1}}$.}}
{\somm Évaluation individuelle de fin de chapitre.\par
 \textit{Indicateurs :} équilibre ($\tan\alpha$) ; potentiel avec la bonne référence ; conservation de l’énergie.}
{45}{1 h 45}
\etape{cGreen}{\axe{cGreen}{ic_reponses}{L’essentiel}{Synthèse du chapitre}{\diapo{42}}}
{\pt{Projette « L’essentiel du chapitre » (diapo 42) ; fait le bilan de l’évaluation ; donne le travail à la maison.}}
{\pt{Relit la synthèse ; note le travail à faire.}}
{\somm Bilan de l’évaluation.}
{15}{2 h 00}
\end{xltabular}

\vspace{2mm}
\noindent
\begin{minipage}[t]{0.49\linewidth}\vspace{0pt}
\encadre{fi_maison}{cOrange}{Travail à la maison et prolongements}{%
\pt{\textbf{Après la séance 1} — refaire l’application 1 avec deux protons d’un noyau distants de \qty{2e-15}{m} : objectif O2.}
\pt{\textbf{Après la séance 2} — exercice 1 (énoncé projeté, diapo 34, à recopier et à résoudre dans le cahier) : objectif O3.}
\pt{\textbf{Après la séance 3} — reprendre l’application 3 pour un proton allant de $P$ vers $N$ : objectif O5.}
\pt{\textbf{Après la séance 4} — exercice 2 (diapo 36) : objectifs O3 et O4.}
\pt{Soutien : représenter les champs de l’application 2 au point situé à \qty{5}{cm} de $B$ hors du segment. Approfondissement : calculer la vitesse d’une particule $\alpha$ de l’exercice 3 ($m = \qty{6,6e-27}{kg}$).}}
\end{minipage}\hfill
\begin{minipage}[t]{0.49\linewidth}\vspace{0pt}
\encadre{ic_application}{cPurple}{Bilan de la séquence (à compléter par le professeur)}{%
{\sffamily\footnotesize Objectifs atteints :}\lignevide\lignevide
\par\vspace{1.5mm}{\sffamily\footnotesize Difficultés rencontrées :}\lignevide\lignevide
\par\vspace{1.5mm}{\sffamily\footnotesize Ajustements pour la prochaine séquence :}\lignevide}
\end{minipage}

\end{document}
